\section{Reporting Failures and Errors}\label{sec:ReportFailure}
An Error may not remain unnoticed, and to some extent, even Failures should be brought to the attention of the users or administrators of our program.

Usually the reporting of an Error is delegated by the aborted operation to its parent operation (in case of a method\index{method} call, the aborted method\index{method} relayed it to the method\index{method} that called it). But at some level, there isn’t a parent operation anymore – so how to report at this level?

\subsection{Ways to Report}
Basically, these are ways an operation (a method\index{method}) has to externalise an error condition (to bring it to the attention of the users or administrators):

\paragraph{Return a status as the return value of the method\index{method} call} – We discussed this already \hyperref[sec:FailureVsError]{above}. It requires that the calling method\index{method} checks the return value to detect the failure, but it does not externalises the error condition. This remains the task of the caller. 

\paragraph{Set a global status flag} – This approach can cause issues in multithreaded environments, because it may be difficult to synchronize concurrent updates to the global field that represents the status flag.

It also requires the calling method\index{method} to check that flag after each call, and it also leaves the externalisation of the error condition to the caller.

You should always prefer a return value for the status over a global status flag, and if that is not feasible, throwing an Exception\index{exception} still might be the better option.

\paragraph{Write a message to the console} – This immediately externalises the error condition, but in the most volatile way. Nevertheless, depending on the context, this may be sufficient.

\paragraph{Open a message window} – Obviously this make sense only for an application with some kind of UI. Usually, such a message window (or \bdq{Message Box}) provides options on how to proceed.

This is similarly volatile as writing the message to the console, but as in most cases the message window requires a user interaction before the program continues, it gets more attention.

\paragraph{Write a message to a file} – This is a simple form of logging, sometimes referred to as \bdq{The Poor Man's Logging}; often this uses the same mechanics as writing a message to a console, just with a different target. It externalises the error condition in a non-volatile way.

\paragraph{Send a message (as email or through a Messaging Service)} – This externalises the error condition in the most consequent manner possible: it sends it into the world. Whether it is volatile or not depends on the receiving entity.

\paragraph{Log a message through a log service} – Logging externalises the error condition by persisting it in some way. How this is done is highly configurable, details about this are discussed in chapter \tqfullvref{sec:Logging}.

\paragraph{Throw an exception\index{exception}} – Signalling an error condition by throwing an exception\index{exception} delegates the whole responsibility to the caller. This will be discussed further in chapter \tqfullvref{sec:ThrowingExceptions}. Depending on the context, it could be useful to define a hierarchy of domain-specific exceptions\index{exception} (refer to chapter \tqfullvref{sec:DomainSpecificExceptions}).

Throwing an exception will not externalise the error condition and it is volatile.

\bigskip

Usually only one way should be used, although it is sometimes a good idea to persist an otherwise volatile status.

The appropriate way to report Failures and Errors depends on the type of product\footnote{as defined in chapter \tqfullvref{sec:TypesOfProducts}} that you work on.

\subsubsection{\nameref{sec:FunctionLibrary}}
A function library\index{product type!function library} should report failures through return values and errors through exceptions\index{exception}. If a library consists of a relatively loose collection of utility functions and classes, a hierarchy of domain-specific exceptions\index{exception} (refer to \tqfullvref{sec:DomainSpecificExceptions}) does not make sense.

Furthermore, the library's methods should not perform any logging; instead, they should delegate this task to the calling methods within the consuming application.

\subsubsection{\nameref{sec:FeatureLibrary}}
A feature library\index{product type!feature library} usually defines a \bdq{domain} and should therefore come a with a hierarchy of domain-specific exceptions\index{exception} that are used to report error conditions to the host application. Other than a function library\index{product type!function library}, the methods\index{method} of a feature library\index{product type!feature library} should never emit standard exceptions to a calling method outside the library.

A feature library\index{product type!feature library} may also log failures and errors through a log services, but it should use a facade for the log service, like SLF4J\index{SLF4J} \autocite{SLF4J}, instead of referring the log service directly.

\subsubsection{\nameref{sec:Tool}\index{product type!tool}}
\subsubsection{\nameref{sec:StandaloneApplication}\index{product type!standalone application}}
\subsubsection{\nameref{sec:ServerbasedApplication}\index{product type!server-based application}}

\subsubsection{\nameref{sec:Extension}}
Extensions\index{product type!extension} are executed in the context of another application that defines how failures and errors must be reported. Usually such an application will provide an API for messaging and/or logging. For the details you have to refer to the documentation of the respective application.

\subsubsection{\nameref{sec:ScriptInJava}}
A script\index{product type!script} should write messages about failures and errors directly to the console.


%==============================================================================
% End of file!!
%==============================================================================

